\section{Introduction}

\subsection{Motivation and research context}

Satoshi Nakamoto published the Bitcoin concept in 2008. And on January 5, 2009, the first ``genesis block'' was generated, and thus the first Bitcoin network came into operation, involving computers to generate Bitcoin blocks. Since then, the technology has evolved from an experiment to a class of high-risk assets. As a result, the cryptocurrency market has now grown so large that it exerts a real influence on the market for conventional assets, as many investors—as well as government funds—use Bitcoin both as a vehicle for storing capital and for conducting financial transactions. This creates a large pool of financial capital, the rise and fall of which can influence conventional, traditional markets. Consequently, the task arises of managing risks and identifying correlations between Bitcoin and traditional markets.

This thesis explores time-varying intermarket dependency. The time-varying statistical relationship between Bitcoin and traditional asset markets:
1) stock indices;
2) gold;
3) silver;
4) the dollar.
Is such a correlation even predictable? Does the cryptocurrency market contain signals about the dynamics of traditional asset prices? To assess market risks and stress, the direction of the market is important, since the cryptocurrency market is more volatile; because of this, it is possible to predict the nature of the movement of traditional asset markets more quickly.

Let's call the investor who needs a tool for precise yield forecasting a ``risk manager.'' A risk manager must be able to identify the moment when the nature of the markets' joint movement changes—the moment of transition from a closely correlated to a decorrelated state, and vice versa. When the likelihood of a stressful day in traditional asset markets increases—given the scale of investments—even a small signal can be very useful, allowing the Risk Manager to withdraw a portion of assets before the markets fall.

\subsection{Research questions}

This thesis is structured around two interrelated levels.
1) The forecasting problem: Is it possible to predict the moving correlation between Bitcoin and a traditional asset (all assets are considered individually; the goal is not to identify an average traditional asset)? Compare machine learning models with classical economic valuation models.
2) Is it possible to turn machine learning forecasts into an early warning signal for real-world investors?

Four research questions are also examined:

\begin{enumerate}
\item Is there a correlation between traditional assets and Bitcoin?
    \item Which is better: the DCC-GARCH method or machine learning models?
    \item Which is better: basic models, such as AR(1), or machine learning models?
    \item Do forecasts based on cryptocurrency data contain practically useful information about stress in traditional markets?
\end{enumerate}

\subsection{Contributions of the thesis}

In this work, a reproducible Python pipeline was designed from scratch and hosted on GitHub. When run, it loads price data from open exchange sources, constructs logarithmic return series, and calculates rolling correlations for various window lengths. The models used include the naive basic random walk model, AR(1) autoregression, the HAR autoregressive model, regularized linear models such as Ridge and ElasticNet, tree ensembles (Random Forest, GBM, XGBoost), and the DCC-GARCH method with data leakage protection. After the first layer completes its work, the second level performs calculations using the first layer's forecasts to determine whether the market is under stress or not.

The main thing that we found out is that we can predict the further development of the addiction process, but most of this predictability is due to what happened in the past. Adding new information from other sources does not bring significant benefits. The Ridge, AR(1), and HAR models eventually turned out to be approximately equally accurate, since they all rely on the fact that the sliding correlation remains very stable over time. Interestingly, these three completely different models came to the same result — this indicates that the objective function has a fairly simple linear structure, and complicating the models does not really add anything. At the same time, machine learning does bring some benefits, especially when it comes to classifying stressful periods. In this case, regularized linear classifiers work better than ensembles of trees.

This study makes three main contributions. First, we developed a reproducible pipeline in Python that covers the entire analysis process from processing the raw data to generating the final plots. Second, we evaluated ten forecasting models, ranging from simple benchmark specifications to more complex approaches such as XGBoost and DCC-GARCH. This provides a consistent framework for assessing whether the added complexity of machine learning methods leads to improved forecasting performance. Third, we used correlation forecasts to build a simple stress scenario classifier, linking the forecasting results to practical applications in the field of risk management.

\subsection{Thesis outline}

Section 2 provides an overview of the theory and literature: time-varying dependence, financial contagion, machine learning-based forecasting methods, and the Dibold–Mariano test. Section 3 describes the data, the construction of the target series, feature engineering, model specifications, and the estimation procedure. Section 4 presents the empirical results for both levels. Section 5 interprets and refines these results. Section 6 summarizes the findings.
